% Dokumentklasse
\documentclass[11pt,titlepage]{report} % Für längere Arbeiten mit Kapiteln und Abschnitten, geeignet für Berichte und Dokumentationen.

% Sprachunterstützung
\usepackage[ngerman]{babel} % Sorgt für sprachspezifische Regeln wie Silbentrennung und deutsche Bezeichnungen.

% Symbole und Icons
\usepackage{fontawesome} % Ermöglicht die Nutzung von Font Awesome Icons.

% Seitenlayout
\usepackage[a4paper,lmargin={3.5cm},rmargin={2.5cm},tmargin={3cm},bmargin={2cm},headsep=0.5cm,headheight=1.5cm]{geometry} % Einstellung von Seitenrändern für A4-Papier.

% Mathematische Pakete
\usepackage{amssymb} % Enthält zusätzliche mathematische Symbole.
\usepackage{amsthm} % Unterstützung für Theoreme, Definitionen und Beweise.
\usepackage{amsmath} % Verbesserte Darstellung von mathematischen Formeln.

% Grafiken und Bilder
\usepackage{graphicx} % Ermöglicht das Einfügen von Bildern und Grafiken.
\usepackage{svg} % Ermöglicht das Einfügen von SVG-Grafiken.

% Typografie
\usepackage{microtype} % Verbessert die Lesbarkeit und Silbentrennung.
\usepackage{fontspec} % Für die Verwendung von TrueType- und OpenType-Schriften.

% Kopf- und Fußzeilen
\usepackage{fancyhdr} % Ermöglicht die Gestaltung von Kopf- und Fußzeilen.

% Farben
\usepackage{xcolor} % Definiert und verwendet Farben im Dokument.

% Weitere mathematische Symbole
\usepackage{amsfonts} % Zusätzliche Schriftarten für mathematische Symbole.

% Tabellen
\usepackage{booktabs} % Verbesserte Gestaltung von Tabellen.
\usepackage{multirow} % Ermöglicht Zellen über mehrere Zeilen in Tabellen.
\usepackage{enumitem} % Anpassung von Listen und Aufzählungen.
\usepackage{tabularx} % Erweiterte Tabellenfunktionen, insbesondere für Spaltenbreiten.

% Literaturverwaltung
\usepackage[notes,backend=biber,autocite=footnote,ibidtracker=false]{biblatex-chicago} % Für Literaturverzeichnisse nach Chicago-Stil (Notes and Bibliography).

% Zitate
\usepackage{csquotes} % Korrekte Darstellung von Zitaten, z. B. Anführungszeichen.

% Fußnoten-Formatierung
\usepackage[hang,flushmargin,bottom,multiple]{footmisc} % Keine Einrückung, mehr Abstand zwischen Fußnoten
\setlength{\footnotesep}{12pt} % Abstand zwischen Fußnoten

% Verhindere Seitenumbrüche in Fußnoten
\interfootnotelinepenalty=10000 % Verhindert Umbrüche innerhalb von Fußnoten
\raggedbottom % Erlaubt variable Seitenhöhen zur besseren Fußnoten-Platzierung

% Zusätzliche Fußnoten-Kontrolle
\setlength{\skip\footins}{10pt plus 4pt minus 2pt} % Flexibler Abstand zwischen Text und Fußnoten
\renewcommand{\footnoterule}{\kern-3pt\hrule width 0.4\columnwidth\kern 2.6pt} % Angepasste Fußnotenlinie

% Zeilenabstand
\usepackage[onehalfspacing]{setspace} % Erzeugt 1,5-fachen Zeilenabstand.

% Hyperlinks
\usepackage[hidelinks]{hyperref}% Fügt klickbare Hyperlinks ein, ohne sie farblich hervorzuheben.

% Positionierung
\usepackage{float} % Präzise Platzierung von Tabellen und Abbildungen.

% Abkürzungen
\usepackage{acronym} % Verwaltung von Abkürzungen

% Glossare
\usepackage[toc=false]{glossaries-extra} % Verwaltung von Glossaren ohne Inhaltsverzeichnis-Eintrag.

% Tabellenfarben
\usepackage{colortbl} % Ermöglicht das Einfärben von Tabellenzellen.

% Tabellen- und Spaltenformatierung
\usepackage{array} % Zusätzliche Optionen für Tabellen und Spalten.

% Diagramme und Visualisierungen
\usepackage{tikz} % Erstellung von Vektorgraphiken und Diagrammen.
\usepackage{pgfplots} % Erstellung von Plots und Grafiken.

\usetikzlibrary{shapes,arrows,positioning,fit,backgrounds,calc,matrix,decorations.pathreplacing,decorations.markings,shapes.geometric,mindmap}

% PDF-Integration
\usepackage{pdfpages} % Einfügen von ganzen PDF-Dokumenten.

% Querformat
\usepackage{pdflscape} % Ermöglicht die Nutzung von Seiten im Querformat.

% Nummerierungsanpassung
\usepackage{chngcntr} % Ermöglicht Änderungen der Nummerierung von Kapiteln und Abschnitten.
\counterwithout{figure}{chapter} % Abbildungen durchgehend nummerieren (ohne Kapitelpräfix)
\counterwithout{footnote}{chapter} % Fußnoten durchgehend nummerieren (ohne Kapitelpräfix)
\counterwithout{table}{chapter} % Tabellen durchgehend nummerieren (ohne Kapitelpräfix)

% Inhaltsverzeichnis
\usepackage{tocloft} % Anpassung des Layouts des Inhaltsverzeichnisses.

% Titelgestaltung
\usepackage{titlesec} % Anpassung von Kapitel- und Abschnittstiteln.

% Ausrichtung
\usepackage{ragged2e} % Ermöglicht den Blocksatz oder Flattersatz.

\pgfplotsset{compat=1.18}

% Pfade für Bilder und andere Ressourcen
\graphicspath{{bilder/}}

% Arial als Hauptschrift (11pt)
\setmainfont{Arial}

% Bibliographie
\addbibresource{quellen.bib}

% Farben
\definecolor{link}{HTML}{0000EE}

% Absatzformatierung
\setlength{\parindent}{0pt}           % Keine Einrückung am Absatzbeginn
\setlength{\parskip}{6pt}             % 6pt Abstand nach jedem Absatz

% Anpassung \listoffigures
\renewcommand{\cftloftitlefont}{\normalsize\bfseries}
\renewcommand{\cftbeforeloftitleskip}{10pt}
\renewcommand{\cftafterloftitleskip}{10pt}

% Anpassung \listoftables
\renewcommand{\cftlottitlefont}{\normalsize\bfseries}
\renewcommand{\cftbeforelottitleskip}{10pt}
\renewcommand{\cftafterlottitleskip}{10pt}

% Anpassung \tableofcontents
\renewcommand{\cfttoctitlefont}{\normalsize\bfseries}
\renewcommand{\cftbeforetoctitleskip}{10pt}
\renewcommand{\cftaftertoctitleskip}{10pt}

% Füge Punkte für Kapitel im Inhaltsverzeichnis hinzu
\renewcommand{\cftchapaftersnum}{.}
\renewcommand{\cftchapdotsep}{\cftdotsep}
\renewcommand{\cftsecaftersnum}{.}

% Präfixe für Abbildungs- und Tabellenverzeichnis
\renewcommand{\cftfigpresnum}{Abb.~}
\renewcommand{\cfttabpresnum}{Tab.~}

% Abstand nach Nummern im Abbildungs- und Tabellenverzeichnis
\setlength{\cftfignumwidth}{4em}
\setlength{\cfttabnumwidth}{4em}

% Überschriftenformatierung
\titleformat{\chapter}
  {\fontsize{11}{16.5}\bfseries\raggedright}
  {\thechapter.}{1em}{}
\titlespacing*{\chapter}{0pt}{0pt}{6pt}

\titleformat{\section}
  {\fontsize{11}{16.5}\bfseries\raggedright}
  {\thesection.}{1em}{}
\titlespacing*{\section}{0pt}{0pt}{6pt}

\titleformat{\subsection}
  {\fontsize{11}{16.5}\bfseries\raggedright}
  {\thesubsection.}{1em}{}
\titlespacing*{\subsection}{0pt}{0pt}{6pt}

\titleformat{\subsubsection}
  {\fontsize{11}{16.5}\bfseries\raggedright}
  {\thesubsubsection.}{1em}{}
\titlespacing*{\subsubsection}{0pt}{0pt}{6pt}

% Blocksatz für Fließtext
\justifying

% Bibliographie-Einstellungen für Chicago-Stil
\DeclareNameAlias{sortname}{family-given}
\DefineBibliographyStrings{german}{andothers = {{et\,al\adddot}},}

% Benutzerdefinierte Zitatbefehle
\newcommand{\customcite}[2][]{%
  \ifblank{#1}{%
    \autocite{#2}%
  }{%
    \autocite[#1]{#2}%
  }%
}

% Fix lange URLs
\newcommand*\oldurlbreaks{}
\let\oldurlbreaks=\UrlBreaks
\renewcommand{\UrlBreaks}{\oldurlbreaks
  \do\a\do\b\do\c\do\d\do\e\do\f\do\g\do\h\do\i\do\j\do\k\do\l\do\m\do\n\do\o\do\p\do\q\do\r\do\s\do\t\do\u\do\v
  %\do\w%
  \do\x\do\y\do\z
  \do?\do&\do_
  \do\0\do\1\do\2\do\3\do\4\do\5\do\6\do\7\do\8\do\9
  \do:\do/
  }
\urlstyle{same}

% Befehle für Links und Verweise
\newcommand{\sct}[2]{\textcolor{link}{\hyperref[#1]{#2}}}
\newcommand{\gl}[1]{\gls{#1}}

% Promptverzeichnis erstellen
\newcommand{\listpromptname}{Promptverzeichnis}
\newlistof{prompt}{prm}{\listpromptname}
% Formatierung des Promptverzeichnis-Titels (wie bei anderen Verzeichnissen)
\renewcommand{\cftprmtitlefont}{\normalsize\bfseries}
\renewcommand{\cftbeforeprmtitleskip}{10pt}
\renewcommand{\cftafterprmtitleskip}{10pt}

% Befehl für Prompt-Referenzen
\newcounter{promptcounter}
\newcommand{\promptentry}[1]{%
  \refstepcounter{promptcounter}%
  \addcontentsline{prm}{prompt}{\protect\numberline{\thepromptcounter}#1}%
  #1%
}

% Inhaltsverzeichnis-Einrückung
\setlength{\cftsubsecindent}{1.5em}
\setlength{\cftsubsecindent}{3.0em}

% Anhangsverzeichnis erstellen
\newcommand{\listappendixname}{Verzeichnis der Anhänge}
\newlistof{appendix}{app}{\listappendixname}

% Formatierung des Anhangsverzeichnis-Titels (wie bei anderen Verzeichnissen)
\renewcommand{\cftapptitlefont}{\normalsize\bfseries}
\renewcommand{\cftbeforeapptitleskip}{10pt}
\renewcommand{\cftafterapptitleskip}{10pt}

% Setze die Breite für Anhangsnummern im Verzeichnis
\newlength{\appnumwidth}
\setlength{\appnumwidth}{6em}

\newcommand{\appendixentry}[1]{%
  \refstepcounter{appendix}%
  \addcontentsline{app}{appendix}{\protect\makebox[\appnumwidth][l]{Anhang \theappendix:}#1}%
}

% Glossar
\makenoidxglossaries
\newglossaryentry{ki}{
  name={KI},
  description={\emph{Künstliche Intelligenz} bezeichnet die Fähigkeit von Maschinen, menschenähnliche, intelligente Verhaltensweisen zu zeigen. Diese Technologie umfasst Methoden wie maschinelles Lernen, neuronale Netze und Expertensysteme.}
}
\glsaddall
% Anpassung der Kopfzeilen
% Behält die Nummer bei nummerierten Kapiteln, entfernt nur das Wort "Kapitel"
\renewcommand{\chaptermark}[1]{\markboth{\thechapter. #1}{}}

% Seitenstil für alle Seiten
\fancypagestyle{plain}{%
  \fancyhf{} % Alle Kopf- und Fußzeilenfelder löschen
  \lhead{\nouppercase{\leftmark}} % Kapitelüberschrift oben links
  \rhead{\thepage} % Seitenzahl oben rechts
}


% Seitenstil für normale Seiten
\pagestyle{fancy}
\fancyhf{} % Alle Kopf- und Fußzeilenfelder löschen
\lhead{\nouppercase{\leftmark}} % Kapitelüberschrift oben links
\rhead{\fontsize{9}{13.5}\selectfont\thepage} % Seitenzahl oben rechts, 9pt
\lhead{\fancyplain{}{\nouppercase{\leftmark}}} % Kapitelüberschrift oben links

% Anpassung der Verzeichnis-Überschriften
\renewcommand{\contentsname}{Inhalt}
\renewcommand{\listfigurename}{Abbildungsverzeichnis}
\renewcommand{\listtablename}{Tabellenverzeichnis}
\renewcommand{\bibname}{Literaturverzeichnis}

% Anpassung der Bezeichnungen für Abbildungen und Tabellen
\addto\captionsngerman{%
  \renewcommand{\figurename}{Abb.}%
  \renewcommand{\tablename}{Tab.}%
}

% Setze PDF-Metadaten
\hypersetup{
  pdftitle={Titel der Arbeit},
  pdfauthor={Vorname Nachname},
  pdfsubject={Fachbereich/Studiengang},
  pdfkeywords={Schlüsselwort1, Schlüsselwort2, Schlüsselwort3},
  pdfcreator={luaLaTeX},
  pdfproducer={pdfLaTeX},
  pdflang={de-DE},
}

\begin{document}

% Schmutztitel (Vorderdeckel)
\begin{titlepage}
\thispagestyle{empty}
\begin{titlepage}
    \thispagestyle{empty}
    \hspace{-1cm}\includegraphics[width=0.55\textwidth]{rh-logo}\\[1.9cm]
    
    \vspace{2cm}
    \begin{minipage}{\textwidth}
    \hspace{2.5cm}\begin{minipage}{\dimexpr\textwidth-2.5cm\relax}
    {\Large\textbf{Fachbereich Ingenieurwesen}}
    
    {\Large Studiengang Informatik (B.Sc.)}\\
    \\
    \\
    
    
    {\Large\textbf{Subjekt}}
    
    \parbox{\textwidth}{\Large\raggedright Titel}\\
    \\
    
    {\normalsize\fontsize{11pt}{16.5pt}\selectfont vorgelegt von}
    
    {\Large Mustermensch}
    
    {\normalsize\fontsize{11pt}{16.5pt}\selectfont Matrikelnr. 111111111}\\
    \\
    
    {\normalsize Sommersemester 2025}
    \end{minipage}
    \end{minipage}
    
    \vfill
    
\end{titlepage}

\end{titlepage}

% Leeres weißes Blatt
\newpage
\thispagestyle{empty}
\quad 
\newpage

% Haupttitel
\begin{titlepage}
\thispagestyle{empty}
\begin{titlepage}
    \thispagestyle{empty}
    \hspace{-1cm}\includegraphics[width=0.55\textwidth]{rh-logo}\\[1.9cm]
    
    \vspace{2cm}
    \begin{minipage}{\textwidth}
    \hspace{2.5cm}\begin{minipage}{\dimexpr\textwidth-2.5cm\relax}
    {\Large\textbf{Fachbereich Ingenieurwesen}}\\
    
    {\Large Studiengang Informatik (B.Sc.)}\\
    \\
    \\
    
    
    {\Large\textbf{Subjekt}}
    
    \parbox{\textwidth}{\Large\raggedright Titel}\\
    \\
    
    {\normalsize\fontsize{11pt}{16.5pt}\selectfont vorgelegt von}
    
    {\Large Mustermensch}
    
    {\normalsize\fontsize{11pt}{16.5pt}\selectfont Matrikelnr. 111111111}\\


    {\normalsize\fontsize{11pt}{16.5pt}\selectfont Erstprüfer: Prof. Dr. Mustermensch}

    {\normalsize\fontsize{11pt}{16.5pt}\selectfont Zweitprüfer: Prof. Dr. Mustermensch}\\
    \\
    
    {\normalsize Sommersemester 2025}
    \end{minipage}
    \end{minipage}
    
    \vfill
    
\end{titlepage}

\end{titlepage}

% Deaktiviere Einrückungen
\setlength\parindent{0pt}

\begin{onehalfspace}
    % Römische Seitennummerierung für Verzeichnisse
    \pagenumbering{roman}
    
    % Zusammenfassung (deutsch) und Abstract (englisch)
    \chapter*{Zusammenfassung}
\addcontentsline{toc}{section}{Zusammenfassung}

% Hier kommt die deutsche Zusammenfassung der Arbeit
Lorem ipsum dolor sit amet, consetetur sadipscing elitr, sed diam nonumy eirmod tempor invidunt ut labore et dolore magna aliquyam erat, sed diam voluptua. At vero eos et accusam et justo duo dolores et ea rebum. Stet clita kasd gubergren, no sea takimata sanctus est Lorem ipsum dolor sit amet.

\vspace{1cm}

\chapter*{Abstract}
\addcontentsline{toc}{section}{Abstract}

% Hier kommt die englische Zusammenfassung der Arbeit
Lorem ipsum dolor sit amet, consetetur sadipscing elitr, sed diam nonumy eirmod tempor invidunt ut labore et dolore magna aliquyam erat, sed diam voluptua. At vero eos et accusam et justo duo dolores et ea rebum. Stet clita kasd gubergren, no sea takimata sanctus est Lorem ipsum dolor sit amet.
    \newpage
    
    % Inhaltsverzeichnis
    \phantomsection
    \tableofcontents
    \newpage

    % Optionale Vorbemerkung
    % Auskommentieren, wenn nicht benötigt
    \chapter*{Vorbemerkung}
\addcontentsline{toc}{section}{Vorbemerkung}

% Hier kann eine Vorbemerkung eingefügt werden, z.B. zum Gendern
In dieser Arbeit wird das generische Maskulinum verwendet, wenn im Mittelpunkt der Aussagen keine spezifischen Personen stehen, sondern die Funktion, die Personen ausüben. Sämtliche biologischen Geschlechter und subjektiv empfundenen Zuordnungen sind damit gleichermaßen gemeint.

% Alternativ kann hier auch ein anderer Hinweis zum Gendern stehen, z.B.:
% In dieser Arbeit werden geschlechtsneutrale Formulierungen verwendet, um alle Geschlechter gleichermaßen anzusprechen.
    \newpage
    
    % Hauptteil mit arabischer Seitennummerierung
    \pagenumbering{arabic}

    % Kapitel einbinden
    \foreach \i in {10,11} {%
      \edef\FileName{abschnitte/\i-kapitel}%

      \IfFileExists{\FileName}{%
        \input{\FileName}
      }
      {%
        %file does not exist
      }
	  }

    % Abkürzungsverzeichnis
    \phantomsection
    \addcontentsline{toc}{chapter}{Abkürzungsverzeichnis}
    \label{sec:Abkuerzungsverzeichnis}
    \chapter*{Abkürzungsverzeichnis}
    \newpage

    % Abbildungsverzeichnis
    \phantomsection
    \addcontentsline{toc}{chapter}{Abbildungsverzeichnis}
    \listoffigures
    \newpage
    
    % Tabellenverzeichnis
    \phantomsection
    \addcontentsline{toc}{chapter}{Tabellenverzeichnis}
    \listoftables
    \newpage
    
    % Promptverzeichnis (für KI-Prompts)
    \phantomsection
    \addcontentsline{toc}{chapter}{Promptverzeichnis}
    \markboth{Promptverzeichnis}{} % Manuelles Setzen der Kopfzeile
    \listofprompt
    \newpage

    % Literaturverzeichnis
    \newpage
    \phantomsection
    \addcontentsline{toc}{chapter}{Literaturverzeichnis}
    \label{chap:Literaturverzeichnis}
    \markboth{Literaturverzeichnis}{} % Manuelles Setzen der Kopfzeile
    \printbibliography[title={Literaturverzeichnis}]
    
    % Anhangsverzeichnis
    \newpage
    \phantomsection
    \addcontentsline{toc}{chapter}{Verzeichnis der Anhänge}
    \listofappendix
    \newpage
    
    % Anhänge
    \newpage
    \appendix
    \section{Anhang}
\label{app:anhang}


    % Eigenständigkeitserklärung
    \newpage
    \chapter*{Eigenständigkeitserklärung}
\addcontentsline{toc}{chapter}{Eigenständigkeitserklärung}
Hiermit bestätige ich, dass ich die vorliegende Arbeit selbstständig verfasst und keine 
anderen Publikationen, Vorlagen und Hilfsmittel (z.B. künstliche Intelligenz) als die 
angegebenen benutzt habe. Alle Teile meiner Arbeit, die wortwörtlich oder dem Sinn nach 
anderen Werken entnommen sind, wurden unter Angabe der Quelle kenntlich gemacht. 
Gleiches gilt für von mir verwendete Internetquellen. Ich versichere, dass ich diese Arbeit 
oder nicht zitierte Teile daraus vorher nicht in einem anderen Prüfungsverfahren eingereicht 
habe. Mir ist bekannt, dass meine Arbeit zum Zwecke eines Plagiatsabgleichs mittels einer 
Plagiatserkennungssoftware auf eine ungekennzeichnete Übernahme von fremdem geistigen 
Eigentum sowie auf die Nutzung von künstlicher Intelligenz zur Texterstellung überprüft 
werden kann. Ich versichere, dass die elektronische Form meiner Arbeit mit der gedruckten 
Version identisch ist.


% \chapter*{Declaration of Originality}
\addcontentsline{toc}{chapter}{Declaration of Originality}
% I hereby confirm that I have independently written this work and have not used any 
% publications, templates, or aids (e.g. artificial intelligence) other than those I have 
% indicated. All parts of my work which have been taken literally or correspondingly from other
% publications have been duly acknowledged. This also applies to Internet sources. I confirm
% that I have not previously submitted this work or any unquoted parts thereof in any other 
% examination procedure. I am aware that my work may be checked for plagiarism by means 
% of plagiarism recognition software, as well as for the use of artificial intelligence for text
% creation, in order to verify the integrity of its written content. I also confirm that the 
% electronic form is identical to the printed version. \\

\vspace{2cm}

\begin{center}
\begin{tabular}{p{0.45\textwidth}p{0.45\textwidth}}
    [Ort], \today & \\
    \hline
    % Place, Date & Signature \\
    Ort, Datum & Unterschrift \\
\end{tabular}
\end{center}
    
    % Lebenslauf (nur für Abschlussarbeiten)
    \newpage
    \chapter*{Lebenslauf}
\addcontentsline{toc}{chapter}{Lebenslauf}

\begin{tabular}{p{0.3\textwidth}p{0.7\textwidth}}
\textbf{Persönliche Daten} & \\
Name: & Max Mustermann \\
Geburtsdatum: & 01.01.1990 \\
Geburtsort: & Musterstadt \\
\\
\textbf{Ausbildung} & \\
09/2022 - heute & Studium an der Rheinischen Hochschule Köln \\
 & Studiengang: Informatik (B.Sc.) \\
09/2019 - 07/2022 & Ausbildung zum Fachinformatiker \\
 & Firma XYZ GmbH, Köln \\
09/2010 - 06/2019 & Allgemeine Hochschulreife \\
 & Muster-Gymnasium, Musterstadt \\
\\
\textbf{Berufliche Erfahrung} & \\
08/2022 - heute & Werkstudent \\
 & Firma ABC GmbH, Köln \\
 & Tätigkeitsbereich: Softwareentwicklung \\
\\
\textbf{Besondere Kenntnisse} & \\
Sprachen: & Deutsch (Muttersprache) \\
 & Englisch (fließend) \\
 & Französisch (Grundkenntnisse) \\
IT-Kenntnisse: & Programmiersprachen: Java, Python, C++ \\
 & Datenbanken: MySQL, PostgreSQL \\
 & Weitere: Git, Docker, Linux \\
\end{tabular}

\end{onehalfspace}
\end{document}
